% ID: SEL-D0004-OB-E03
% ESTADO: cotejada. Cotejo literal independiente y revisión de composición realizados.
% Fuente: Selectividad/sel_2019_matematicas_aplicadas/Suplente Septiembre/Exámenes Matemáticas Aplicadas CCSS - A y B.pdf
% SHA-256: eecad753694b3d8cb2d427efddb00d75faae9c4ec5b027056291bc95d7d2e3da
% Página: 2; opción B; ejercicio 3.
% Transcripción de primera pasada basada en la imagen de la página original.
% Se conservan los umbrales de edad y las condiciones de cada apartado.
% Sin solución ni modificación de datos o condiciones.

El 70\,\% de los taxistas de una ciudad tiene 40 años o más y de estos, el 60\,\% es propietario de la licencia del vehículo. Sin embargo, en el caso de los menores de 40 años, son propietarios de la licencia el 23\,\%. Se escoge al azar un taxista de esa ciudad.
\begin{enumerate}[label=\alph*),leftmargin=2.4em]
  \item \textbf{(1 punto)} Calcule la probabilidad de que sea propietario de la licencia del vehículo.
  \item \textbf{(1 punto)} Sabiendo que no es propietario de la licencia, calcule la probabilidad de que tenga 40 años o más.
  \item \textbf{(0.5 puntos)} Calcule la probabilidad de que sea propietario de la licencia o tenga menos de 40 años.
\end{enumerate}
